% Generated draft: requires mathematical and source review.
\documentclass[11pt]{article}
\usepackage[margin=25mm]{geometry}
\usepackage{fontspec}
\setmainfont{Latin Modern Roman}
\usepackage{amsmath,amssymb,mathtools,unicode-math}
\setmathfont{Latin Modern Math}
\usepackage{enumitem,needspace,xurl,hyperref}
\setlength{\parindent}{0pt}
\setlength{\parskip}{5pt}
\setlength{\emergencystretch}{3em}
\newcommand{\N}{\mathbb{N}}
\newcommand{\Z}{\mathbb{Z}}
\newcommand{\Q}{\mathbb{Q}}
\newcommand{\R}{\mathbb{R}}
\newcommand{\C}{\mathbb{C}}
\newcommand{\F}{\mathbb{F}}
\newcommand{\A}{\mathbb{A}}
\newcommand{\PP}{\mathbb{P}}
\newcommand{\E}{\mathbb{E}}
\newcommand{\eps}{\varepsilon}
\newcommand{\dd}{\,\mathrm{d}}
\begin{document}
\section*{UnsolvedMath ID: 30005522}
\textbf{Source problem code:} catalogue entry
\subsection{Definitions and mathematical statement}
The supplied problem asks for a classification of functions whose tomographic data permit unique recovery of a continuous motion. The admissible functions, the forward measurement operator, the allowed motions, and the equivalence under which recovery is unique are not specified in the supplied record. Consequently, a fully determined mathematical statement cannot be recovered from this record alone.

A conditional formulation can make the missing specifications explicit. Let \(\mathcal A\) denote an as-yet unspecified class of admissible functions, let \(\mathcal M\) denote an as-yet unspecified class of continuous motions, and let \(\mathcal D(f,m)\) denote the complete tomographic data generated by \(f\in\mathcal A\) undergoing \(m\in\mathcal M\). Let \(\sim\) denote the equivalence of motions allowed by the intended notion of recovery. If the function \(f\) is known and uniqueness is required for every admissible motion, the classification target would be
\[
\mathcal U=\left\{f\in\mathcal A:\ \forall m_1,m_2\in\mathcal M,\quad \mathcal D(f,m_1)=\mathcal D(f,m_2)\ \Longrightarrow\ m_1\sim m_2\right\}.
\]
The task under those specifications would be to characterize \(\mathcal U\) by necessary and sufficient conditions on \(f\). This is a schematic formulation only: the source does not determine these classes and operators, whether \(f\) is known, whether uniqueness is required for every motion, or whether \(\sim\) means exact equality or equality modulo specified ambiguities. If \(f\) is also unknown, the comparison must instead allow competing functions, and the displayed condition does not express that joint inverse problem.

The accompanying extracted theorem fragment assumes a rotation
\[
R\in C^1([0,T];\mathrm{SO}(3)),\qquad
\mathrm{SO}(3)=\{Q\in\R^{3\times3}:Q^{\mathsf T}Q=I_3,\ \det Q=1\},
\]
and an associated angular velocity \(\omega:[0,T]\to\R^3\). Its claimed relation is stated for every \(r\in(-k_0,k_0)\). However, the measurement formula and differential identity are corrupted, and their remaining symbols are undefined. These fragmentary hypotheses do not establish the motion class or measurement model of the classification problem. In particular, the source's use of continuous motion must not silently be replaced by continuously differentiable rotation.
\subsection{Short English statement}
Classify the functions that allow a continuous motion to be determined uniquely from tomographic measurements. The supplied record leaves the measurement model, function class, motion class, and meaning of uniqueness unresolved, so the displayed mathematical formulation is conditional and requires human review.
\subsection{Sources}
\begin{itemize}
\item Local UnsolvedMath registry, \texttt{data/unsolvedmath/problems.json}, permanent ID \(30005522\), title \emph{Unique Recovery of Continuous Motions}, exact source problem code \texttt{catalogue entry}. The record gives proposed year 2023 and identifies no proposer. Its labels \texttt{open} and \texttt{corrected\_verified}, and its assertion that the official report was checked directly, are supplied registry metadata, not independent verification in this draft.
\item Supplied report citation: \emph{Tomographic Inverse Problems: Mathematical Challenges and Novel Applications} (2023), Oberwolfach Reports, pp. 1105--1194, DOI 10.4171/owr/2023/21. \url{https://doi.org/10.4171/owr/2023/21}. The registry describes this collection as arising from workshop open-problem discussions.
\item Supplied literature source labelled ``common-circle primary paper'': \url{https://arxiv.org/abs/2209.08022}. The record supplies no authors or title for this reference.
\item Supplied literature assessment, checked on exactly \texttt{2026-08-22}: common-circle methods provide uniqueness and motion-recovery results under hypotheses, while the classification for continuous motion is labelled open. The listed sources are the preceding arXiv paper and Oberwolfach report. This is supplied dated curation, described by the record as literature triage; its status and cited claims have not been independently verified here.
\item The extracted passage refers to method-of-moments results as reference [7], common-circle results as reference [6], and a theorem as reference [8]. Their bibliographic identities and correspondence to the supplied links are not established by the record.
\end{itemize}
\end{document}
